\section{Line Integrals}
\label{sec:line-integrals}
\subsection{Introduction---From Local Quantities to Global Physics}
\label{sec:local-to-global-physics}

\subsection{The transition from local to global descriptions}

In the previous chapters we developed the mathematical language required to
describe quantities at individual points in space. We introduced vectors,
scalar and vector fields, coordinate systems, and the differential operators
gradient, divergence, and curl. These concepts allow us to answer questions
such as

\begin{itemize}
  \item What is the electric field at this point?
  \item What is the temperature here?
  \item How rapidly is the pressure changing?
  \item Is the fluid rotating locally?
\end{itemize}

These are \emph{local} questions. They concern the behavior of physical
quantities in an infinitesimally small neighborhood of a single point.

Many of the most important physical questions, however, are not local. They
concern the accumulated behavior of a field over an extended region of space.
For example:

\begin{itemize}
  \item How much work does a force perform while moving an object along a path?
  \item How much electric charge is enclosed inside a surface?
  \item How much fluid passes through a pipe every second?
  \item How much mass is contained inside a planet?
  \item What is the total magnetic flux through a superconducting loop?
\end{itemize}

These questions cannot be answered by differentiation alone. They require
integration.

\subsection{From infinitesimal pieces to complete systems}

Integration is based on a remarkably simple idea. Suppose we wish to determine
the length of a winding mountain road. Instead of measuring the entire road at
once, we divide it into many extremely small segments. Each segment is nearly
straight and therefore easy to measure. If the scalar length element is denoted
by $ds$, then adding all infinitesimal pieces gives the total length

\begin{equation}
  \boxed{L=\int_C ds.}
\end{equation}

The symbol $C$ denotes the curve followed by the road. It is important to
distinguish the vector displacement $d\mathbf r$ from its scalar magnitude:

\begin{equation}
  ds=\lvert d\mathbf r\rvert.
\end{equation}

The same principle extends far beyond geometry. Instead of adding small
lengths, we may add infinitesimal amounts of work, area, volume, electric
charge, probability, or energy. Integration is therefore the mathematical
operation that converts local information into global quantities.

\subsection{Differential and integral descriptions}

Throughout theoretical physics there exist two complementary descriptions of
nature. The first describes what happens at each point. The second describes
what happens over an entire region.

\begin{center}
\begin{tabular}{ll}
\toprule
\textbf{Local description} & \textbf{Global description} \\
\midrule
Differential equations & Integral equations \\
Pointwise behavior & Total accumulated behavior \\
Derivatives & Integrals \\
Local geometry & Global geometry \\
Infinitesimal change & Finite quantity \\
\bottomrule
\end{tabular}
\end{center}

Neither description is universally more fundamental. Under suitable
smoothness and geometric conditions, they are equivalent formulations of the
same physical law. One of the deepest achievements of mathematical physics is
the demonstration of this equivalence.

\subsection{Work performed by a variable force}

In elementary mechanics the expression

\begin{equation}
  W=\mathbf F\cdot\mathbf d
\end{equation}

gives the work done by a constant force during a fixed displacement
$\mathbf d$. Suppose instead that the force varies continuously from one point
to another, as it does in gravitational and electric fields, springs, and many
forms of drag.

We divide the path into infinitesimal displacements $d\mathbf r$. Over each tiny
segment the force is approximately constant, so the infinitesimal work is

\begin{equation}
  dW=\mathbf F\cdot d\mathbf r.
\end{equation}

Adding every contribution along the path $C$ gives

\begin{equation}
  \boxed{W=\int_C \mathbf F\cdot d\mathbf r.}
\end{equation}

This expression is a \emph{line integral}. The dot product selects the component
of the force tangent to the motion. A force perpendicular to the instantaneous
displacement contributes no work at that point.

\subsection{Flux through a surface}

Consider rain falling onto a roof. Instead of asking how strong the rain is at
one point, we ask how much rain reaches the entire roof each second. The answer
depends on the local rain field, the area of the roof, and the orientation of
each piece of the surface.

For an infinitesimal surface element of area $dA$ with unit normal
$\hat{\mathbf n}$, define the oriented area vector

\begin{equation}
  d\mathbf A=\hat{\mathbf n}\,dA.
\end{equation}

The total flux of a vector field $\mathbf F$ through a surface $S$ is

\begin{equation}
  \boxed{\Phi=\int_S \mathbf F\cdot d\mathbf A.}
\end{equation}

Flux appears throughout physics in electric and magnetic fields, fluid flow,
heat transport, gravitational fields, and quantum probability current.

\subsection{Quantities distributed throughout a volume}

Some physical quantities occupy an entire three-dimensional region. Suppose a
material has mass density $\rho(x,y,z)$. The density at a single point does not
determine the total mass. We divide the object into infinitesimal volume
elements $dV$. The mass of one element is

\begin{equation}
  dm=\rho\,dV.
\end{equation}

Adding all contributions yields

\begin{equation}
  \boxed{M=\int_V \rho\,dV.}
\end{equation}

The same construction applies to electric charge, particle number, internal
energy, entropy, and quantum probability.

\subsection{Why the integral theorems matter}

At first glance, differentiation and integration appear to be different
operations. The great integral theorems of vector calculus reveal that they are
deeply connected. In schematic form,

\begin{equation}
  \boxed{\text{boundary information}\quad\longleftrightarrow\quad
  \text{interior information}.}
\end{equation}

More precisely, circulation around a closed boundary is related to the curl
inside the bounded surface, while flux through a closed surface is related to
the divergence inside the enclosed volume. These relations allow difficult
volume calculations to be replaced by surface calculations, and conversely.
They also provide the mathematical structure behind many conservation laws.

\subsection{Historical perspective}

The development of integral calculus began with Isaac Newton and Gottfried
Wilhelm Leibniz in the late seventeenth century. Their independent formulations
provided systematic methods for studying continuously varying quantities.
During the eighteenth and nineteenth centuries, Leonhard Euler, Joseph-Louis
Lagrange, Carl Friedrich Gauss, George Green, William Thomson, and George
Gabriel Stokes extended these ideas into multidimensional geometry.

James Clerk Maxwell recognized that these mathematical results were ideally
suited to electric and magnetic fields. In Maxwell's theory, integral theorems
are not mathematical curiosities: they connect local field equations to
measurable quantities such as enclosed charge, circulation, and flux.

\subsection{Applications throughout physics}

In classical mechanics, line integrals describe the work performed by variable
forces. In electromagnetism, surface integrals define electric and magnetic
flux, while volume integrals determine total charge. In fluid mechanics, flux
integrals measure transport through boundaries, and divergence theorems express
conservation of mass.

In general relativity, integration over curved spacetime regions is essential
for defining invariant and conserved quantities. In quantum mechanics, volume
integrals normalize the wavefunction:

\begin{equation}
  \boxed{\int_{\mathbb R^3}\lvert\psi(\mathbf r,t)\rvert^2\,dV=1.}
\end{equation}

The probability that a particle lies inside a region $V$ is

\begin{equation}
  P(V)=\int_V \lvert\psi\rvert^2\,dV,
\end{equation}

while probability flow through a surface is expressed by a flux integral of
the probability-current density $\mathbf j$.

\begin{summarybox}
Differential operators describe the local behavior of fields. Integrals add
local contributions over curves, surfaces, and volumes to produce global
physical quantities. Line integrals describe work and circulation, surface
integrals describe flux, and volume integrals convert densities into totals.
The integral theorems connect these global quantities to local divergence and
curl.
\end{summarybox}

\begin{lookingaheadbox}
The next section begins with the geometry of parameterized curves and arc
length. This provides the rigorous foundation for line integrals, work,
circulation, conservative forces, and potential energy.
\end{lookingaheadbox}

\subsection{Definition}
For a curve $C$ parametrized by $\mathbf r(t)$, the line integral of a vector field is
\[
\int_C \mathbf F\cdot d\mathbf r
=\int_a^b \mathbf F(\mathbf r(t))\cdot \mathbf r'(t)\,dt.
\]
It measures accumulated tangential action, such as work done by a force along a trajectory.
